\section{5. Přednáška}

\subsection{Binární vyhledávací strom}

\begin{definition}
Strom nazveme \textbf{binární}, pokud
    \begin{itemize}
        \item je zakořeněný
        \item každý vrchol má nejvýše dva syny
        \item u synů rozlišujeme, který je levý a který pravý
    \end{itemize}
\end{definition} \bigskip

\begin{remark}
Zakořeněnění je důležitá vlastnost. Řekněme, že máme popsat jak ve stromu najít vrchol $y$. Popis se
přeci bude lišit, pokud začneme od kořene $u$ nebo $v$. Kořen ve vyhledávacích stromech budeme
využívat jako jakýsi referenční bod.
\end{remark} \bigskip

\begin{definition}[Značení]
Pro vrchol $v$ binárního stromu $T$ značíme
\begin{itemize}
    \item $\ell(v)$ a $r(v)$: levý a pravý syn vrcholu $v$,
    \item $L(v)$ a $R(v)$: levý a pravý podstrom vrcholu $v$,
    \item $p(v)$: otec vrcholu $v$,
    \item $T(v)$: podstrom tvořený kořenem $v$ a všemi jeho potomky,
    \item $h(T(v))$: hloubka stromu $T(v)$ je počet hladin $T(v)$,
    \item $|T|$, $|T(v)|$, $|L(v)|$ a $|R(v)|$: počet vrcholů stromu $T$, $T(v)$, $L(v)$ a $R(v)$.
\end{itemize}
Pokud vrchol $v$ nemá levého syna, položíme $\ell(v) = \emptyset$. Podobně pro $r(v)$ a $p(v)$.\newline
Pak dodefinujeme, že $T(\emptyset)$ je prázdný strom a $h(T(\emptyset)) = 0$.
\end{definition} \bigskip

\noindent Nyní si zavedeme binární vyhledávací strom. Je důležité tyto pojmy rozlišovat.

\begin{definition}
    Binární vyhledávací strom (BVS) je binární strom, v jehož každém vrcholu v je uložen \textbf{unikátní} klíč k(v) a pro každý jeho vrchol platí:

    \begin{itemize}
        \item Pokud $a \in L(v)$, pak $k(a) \bm{<} k(v)$
        \item Pokud $b \in R(v)$, pak $k(b) \bm{<} k(v)$
    \end{itemize}
\end{definition} \bigskip

\begin{remark}
  Pozorný čtenář si jistě všiml využití ostré nerovnosti. To nám implikuje, že v BVS může být každý
  prvek nanejvýš jednou. Sami si zkuste sestrojit BVS, ve kterým je např. $5$ dvakrát. Narozdíl od
  toho vzpomeňme si na binární haldu, kde duplikáty být mohou díky nerovnosti neostré.
\end{remark}

\newpage

\subsection{BVS Algoritmy}

\begin{remark}
Nyní si ukážeme pár zajímavých operací nad BVS. Tyto operace je lepší zkusit pochopit, než se učit
pseudokód nazpaměť.
\end{remark} \bigskip

\subsubsection{BVSShow}

\noindent První je algoritmus BVSShow. Jeho cílem je provést tzv. InOrder průchod stromu a vypsat všechny klíče vzestupně. Z definice BVS víme, že pro každý vrchol $v$ platí: klíče v jeho levém podstromě jsou menší a klíče v pravém podstromě větší. Aby algoritmus vypsal klíče ve vzestupném pořadí, musí se nejprve co nejvíce zanořit doleva, k nejmenším klíčům. Teprve po návratu z levého syna vypíše klíč aktuálního vrcholu a poté pokračuje do pravého syna, kde jsou klíče větší. Tento postup se opakuje pro každý vrchol stromu.
Jelikož v každém vrcholu strávíme konstantní čas, celková časová složitost algoritmu je $\bigO{|T(v)|}$. \bigskip

\noindent \textbf{Vstup}: Ukazatel na kořen $v$ libovolného BVS $T(v)$ \newline
\noindent \textbf{Výstup}: Vzestupně uspořádaná posloupnost klíčů všech vrcholů v $T(v)$ \newline
\funkce{BVSShow}{ukazatel na vrchol $v$}
\begin{pseudocode}
  \item Pokud $v = \emptyset$: return \textcolor{gray}{//$T(v)$ je prázdný}
  \item Zavoláme \textcolor{blue}{BVSShow}($\ell(v)$)
  \item Vypíšeme $k(v)$
  \item Zavoláme \textcolor{blue}{BVSShow}($r(v)$)
\end{pseudocode} \bigskip

\subsubsection{BVSMin}

\noindent BVSMin/BVSMax. Jeho cílem je vypsat minimální/maximální hodnotu klíče ve stromě. Pro minimální klíč se z daného vrcholu zanořujeme stále doleva. Jakmile vrchol již levého syna nemá, pak je minimálním. Analogicky pro maximum, pouze se zanořujeme doprava. V každé hladině navštívíme jeden vrchol a na každém vrcholu provedeme konstantě operací tedy celkem provedeme $\bigO{h(T(v))}$\bigskip

\noindent \textbf{Vstup}: Ukazatel na vrchol $v$ libovolného BVS $T(v)$ \newline
\noindent \textbf{Výstup}: Minimimální klíč daného BVS $T(v)$ \newline
\funkce{BVSMin}{ukazatel na vrchol $v$}
\begin{pseudocode}
  \item Pokud $v = \emptyset$: return \textcolor{gray}{//$T(v)$ je prázdný}
  \item Pokud $\ell(v) = \emptyset$: return $k(v)$
  \item Zavoláme \textcolor{blue}{BVSMin}($\ell(v)$)
\end{pseudocode} \bigskip

\subsubsection{BVSFind}

\noindent BVSFind slouží k nalezení vrcholu ke klíči v BVS. Dostaneme na vstupu ukazatel na kořen $v$ BVS a hledaný klíč. Nyní postupuje následovně, pokud je hledaný klíč větší než $k(v)$. Tak musíme jít do pravého podstromu $R(v)$, kde jsou klíče větší. Analogicky pro menší klíč. Dokud buďto nenalezneme vrchol s hledaným klíčem nebo nevylezeme ze stromu. Stejně jako u BVSMin v kažé hladině navštívíme jeden vrchol a v každém vrcholu provedeme pouze konstantě mnoho operací a celkem tedy dostaneme $\bigO{h(T(v))}$. \bigskip

\noindent \textbf{Vstup}: Ukazatel na vrchol $v$ libovolného BVS $T(v)$ a hledaný klíč $x$. \newline
\noindent \textbf{Výstup}: Ukazatel na vrchol s klíčem $x$, pokud takový v $T(v)$
existuje, jinak $\emptyset$. \newline
\funkce{BVSFind}{ukazatel na vrchol $v$ a hledaný klíč $x$}
\begin{pseudocode}
    \item Pokud $v = \emptyset$: return $\emptyset$
    \item Pokud $k(v) > x$: return \textcolor{blue}{BVSFind}($\ell(v)$)
    \item Pokud $k(v) < x$: return \textcolor{blue}{BVSFind}($r(v)$)
    \item return $v$
\end{pseudocode} \bigskip

\subsubsection{BVSPred}

\noindent BVSPred a BVSSucc slouží k získání předchůdce/následníka nějakého vrcholu $v$ v BVS $T$. Pro předchůdce, je to vždy takový klíč $x$, který je největší ze všech klíčů menších než $x$ v $T$. Jak ho nalezneme? Rozlišujeme zde dva případy:
\begin{itemize}
    \item Ten lehčí je když vrchol $v$ má levého syna. V tom případě je předchůdce vrchol s maximálním klíčem z celého levého podstromu $L(v)$. To platí, jelikož jsou všechny klíče v $L(v)$ menší než $k(v)$ a maximum z $L(v)$ je ten největší menší.
    \item Druhý případ je když $v$ nemá levého syna v tom případě se musíme vynořit nahoru (doprava jít nemůžeme tam jsou všechny klíče větší). Takto cestujeme směrem ke kořeni až do doby než se vrátíme do rodiče $p$ z pravého syna. Tj. $v = r(p(v))$. Předchůdce je poté vrchol $p$. Zprava se chceme vrátit jelikož my potřebujeme najít první menší klíč. Pokud by jsme přišli z levého syna, tak je rodič větší než celý náš podstrom.
\end{itemize}
Pokud náš prvek je minimálním v celém stromě vrátíme nic. Alternativně lze dodefinovat vrcholy s klíčy $\pm\infty$.\newline
Analogicky platí i pro následovníka. Stačí prohodit slova levý a pravý.
buďto využijeme BVSMax nebo procházíme po rodičích a případně do kořene. celkem tedy za $\bigO{h(T(v))}$. \bigskip

\noindent \textbf{Vstup}: Ukazatel na kořen $v$ libovolného BVS $T(v)$ a nějaký jeho vrchol $w$. \newline
\noindent \textbf{Výstup}: Ukazatel na předchůdce vrcholu $w$ v BVS $T$.
existuje-li, jinak $\emptyset$. \newline
\funkce{BVSPred}{ukazatel na kořen $v$ a vrchol $w$}
\begin{pseudocode}
    \item Pokud $l(w) \neq \emptyset$
    \item \tab return \textcolor{blue}{BVSMax}($l(w)$) \textcolor{gray}{//Máme levého syna, tedy přechůdce je maximem $L(w)$}
    \item $z := p(w)$
    \item Dokud ($z \neq \emptyset$ $\And$ $w = l(z)$) \textcolor{gray}{//Jdeme směrem ke kořeni a čekáme až přijdeme zprava}
    \item \tab $w := z$
    \item \tab $z := p(z)$
    \item return $z$
\end{pseudocode} \bigskip

\subsubsection{BVSInsert}

\noindent BVSInsert jak z názvu plyne vkládá prvek do BVS. Na vstupu dostaneme ukazatel na kořen $v$ nějakého BVS a klíč $x$, který chceme vložit. Ve stromě předpokládáme unikátní klíče v případě již existujícího stejného klíče se vložení neprovede. Algoritmus prochází strom stejně jako BVSFind a hledá volné místo kam vložit prvek. Jakmile se dostaneme mimo strom, tak jsme nalezli místo kam vložit daný klíč. Pokud je strom prázdný, tak pouze vložíme prvek. V každém vrcholu provedeme pouze konstantě mnoho operací a jdeme skrze vrcholy od kořene až do listu, což je celkem za $\bigO{h(T(v))}$. \bigskip

\noindent \textbf{Vstup}: Ukazatel na kořen $v$ libovolného BVS $T(v)$ a vkládaný klíč $x$. \newline
\noindent \textbf{Výstup}: Ukazatel na kořen BVS s klíčem $x$. \newline
\funkce{BVSInsert}{Ukazatel na kořen $v$ BVS a vkládaný klíč $x$}
\begin{pseudocode}
    \item Pokud $v = \emptyset$: \textcolor{gray}{//Strom je prázdný}
    \item \tab Vytvoř nový vrchol $v$
    \item \tab return $v$in overleafin overleaf
    \item Pokud $k(v) > x$:
    \item \tab $l(v) :=$ \textcolor{blue}{BVSInsert}$(l(v), x)$
    \item Pokud $k(v) < x$:
    \item \tab $r(v) :=$ \textcolor{blue}{BVSInsert}$(r(v), x)$
    \item return $v$
\end{pseudocode} \bigskip

\subsubsection{BVSDelete}

\noindent Poslední operací je BVSDelete, ten slouží k odstranění vrcholu $v$ z BVS dle zadaného klíče. Zde máme 3 možnosti:
\begin{itemize}
    \item První je pokud vrchol $v$ nemá žádného syna neboli je listem. V tom případě ho prostě odstraníme a omezení BVS ani souvislost jistě neporušíme.

    \item Další případ je kdy má právě jednoho syna, nazvěme ho $w$. Ani zde není problém, rodiči vrcholu $v$, $p(v)$ nahradíme syna $v$ za vrchol $w$. Vrchol $v$ poté bezproblému odstraníme.

    \item Posledním případem je kdy vrchol $v$ má oba syny. V tomto případě nemůžeme vrchol $v$ odstranit jinak by jsme strukturu rozbili. Pomůžeme si trikem a to pomocí nahrazení vrcholu $v$ za jeho následníka nebo předchůdce (ti musí existovat, jelikož má vrchol $v$ oba syny a zároveň to jsou minima/maxima pravých/levých podstromů) a následným odevráním klíče v $T(v)$. Také jsme nemohli rozbít definici BVS, jelikož jsme nahradili bezprostředním následníkem/předchůdcem a stále jsou vlevo menší klíče a vpravo větší. Zároveň následník vždy nemá levého syna, tedy při dalším volání už je oprava struktury BVS opět jednoduchá (Analogicky předchůdce nemá pravého syna).
\end{itemize}
Hledáme vrchol jako při BVSFind, k tomu hledáme minimum podstromu pomocí BVSMin, opět v každé hladině provedeme konstantě mnoho operací tedy celkem $\bigO{h(T(v)}$. \bigskip

\noindent \textbf{Vstup}: Ukazatel na kořen libovolného BVS $T(v)$ a klíč $x$. \newline
\noindent \textbf{Výstup}: Ukazatel na kořen BVS, ze kterého byl odstraněn vrchol s
klíčem $x$, existoval-li. \newline
\funkce{BVSDelete}{Ukazatel na kořen $v$ BVS a mazaný klíč $x$}
\begin{pseudocode}
    \item Pokud $v = \emptyset$: return $\emptyset$
    \item Pokud $k(v) > x$: 
    \item \tab $l(v) := $ \textcolor{blue}{BVSDelete}($\ell(v), x$)
    \item Pokud $k(v) < x$:
    \item \tab $r(v) := $ \textcolor{blue}{BVSDelete}($r(v), x$)
    \item Pokud $k(v) = x$:
    \item \tab Pokud $\ell(v) = r(v) = \emptyset$: return $\emptyset$ \textcolor{gray}{//Vrchol $v$ je listem, mazání je triviální}
    \item \tab Pokud $\ell(v) = \emptyset$: return $r(v)$ \textcolor{gray}{//Vrchol v má pouze jednoho syna, mazání je triviální}
    \item \tab Pokud $r(v) = \emptyset$: return $\ell(v)$
    \item \tab $w :=$ \textcolor{blue}{BVSMin}($r(v)$)
    \item \tab $k(V) := k(w)$
    \item \tab $r(v) := $ \textcolor{blue}{BVSDelete}($r(v), k(w)$) \textcolor{gray}{//Musíme smazat ještě ten prohozený vrchol}
    \item return $v$
\end{pseudocode} \bigskip

\subsection{Časová složitost operací nad BVS}

\begin{remark}
BVS nemá jednoznačnou strukturu. Skutečne zkuste provést BVSInsert na prvcích \{2, 1, 3\} a poté na \{1, 2, 3\}. Množiny jsou v obou případech stejné. Ale výsledný BVS je jiný.\bigskip
\end{remark}

\noindent Časová složitost a problém BVS. Již jsme určili že operace jako BVSInsert, BVSDelete, BVSMin, BVSFind, atd. mají časovou složitost $\bigO{h(T(v)}$. Zde ale nastává problém, pokud nám strom zdegraduje na spojový seznam. To se stane např. pokud všechny prvky budou mít pouze pravého syna. Poté je $\bigO{h(T(v)} = \bigO{|T(v)|} = \bigO{n}$. Což se nám nelíbí, jelikož ztrácíme výhody oproti běznému poli. Jistě lze namítnout, že šance na tento scénář je malá. V případě kdy, ale přijde útočník a začne nám dávat taková data, která toto způsobí. tak jsme nahraní. Na ochranu proti takovémuto zdegenerování se zavádí pojmy vyváženosti. \bigskip

\subsection{Dokonale Vyvážený BVS}
\begin{definition}
BVS nazveme \textbf{dokonale vyvážený}, pokud pro každý jeho vrchol $v$ platí:
\[ \left| |L(v)| - |R(v)| \right| \leq 1 \]
\end{definition} \bigskip

\begin{remark}
Dokonale vyvážený BVS o velikosti $n$ má $1 + \floor{\log_2 n}$ hladin. Operace na takovém stromě mají tedy optimální časovou složitost $\bigO{\log n}$.
\end{remark} \bigskip

\noindent Dokonalé vyvážení nám krásně zajišťuje logaritmicou hloubku BVS. Bohužel se ukazuje, že udržet tuto strukturu stromu současně s přidáváním a mazáním prvků, je příliš nákladné. Alespoň jedna z operací BVSInsert a BVSDelete zdegradují na lineární složitost $\bigO{n}$. \medskip

\begin{theorem}[O složitosti operací nad dokonale vyváženým BVS]
Pokud má BVS zůstávat po provedení operací BVSInsert a BVSDelete dokonale vyvážený, potom ať použijeme jakoukoli implementaci zmíněných operací, bude časová složitost aspoň jedné z nich $\bigOmega{n}$.
\end{theorem} \bigskip

\noindent\textbf{Myšlenka důkazu:} Stačí nám ukázat že existuje dokonale vyvážený BVS a nějaká posloupnost operací BVSInsert a BVSDelete kde alespoň jedna bude mít složitost $\bigOmega{n}$. \bigskip

\begin{proof}
    Použíjeme důkaz přímý
    \begin{enumerate}
    \item Uvažujme dokonale vyvážený BVS s klíči 1,2,...,n, kde n = $2^k - 1$.
    \item Všimneme si, že celý tento strom je určen jednoznačně a všechna lichá čísla jsou v listech tohoto stromu.
    \item Nyní provedeme dvojice operací BVSInsert(n + 1) a BVSDelete(1), pak BVSInsert(n + 2) a BVSDelete(2), ...
    \item Po provedení $i$-té dvojice operací bude strom obsahovat hodnoty $i + 1,..., i$ + n. V závislosti na tom jestli je $i$ liché nebo sudé se vždy listech budou nacházet sudé nebo liché hodnoty. To ale znamená, že všem vrcholům se změní zda jsou listem nebo nejsou. Proto musíme upravit alespoň $\bigOmega{n}$ ukazatelů. Tedy alespoň jedna z operací BVSInsert nebo BVSDelete trvá $\bigOmega{n}$.
    \end{enumerate}
\end{proof}

% ============================================================
% AVL Strom
% ============================================================

\subsection{AVL strom}

\noindent U úplně vyváženého BVS jsme kvůli přílišnému omezení narazili na degradaci rychlostí operací BVSInsert a BVSDelete. Proto musíme zvolit jinou cestu a to pomocí hloubkového vyvážení. \bigskip

\begin{definition}
    BVS je \textbf{hloubkově vyvážený}, pokud pro každý jeho vrchol $v$ platí 
    \[ \left| h(L(v)) - h(R(v)) \right| \leq 1 \]
\end{definition} \bigskip

\noindent Řešení jak toto provést je více (např. Červeno-Černé stromy viz. BI-DAS). Zde v AG1 si představíme AVL stromy. Prvně ale musíme ukázat, že opravdu toto řešení dosáhne logaritmické hloubky stromu, jinak by nemělo smysl používat. \bigskip

\subsubsection{Důkaz hloubky AVL}

\begin{theorem}[o hloubce hloubkově vyvážených stromů]
    Hloubkově vyvážený strom s $n$ vrcholy má hloubku $\bigTheta{log(n)}$.
\end{theorem} \medskip

\noindent \textbf{Myšlenka důkazu:} ukážeme že nejmenší možný strom (dle počtů vrcholů) stále zachovává logaritmickou hloubku což nám dá $\bigOmega{\log_2{(n)}}$. Poté díky tomu, že je binární víme, že pro největší možný strom platí $\bigO{\log_2{(n)}}$. Víme z důkazu o hloubce binárního stromu. Tedy celkem $\bigTheta{\log_2{(n)}}$. \medskip

\begin{proof}
    Důkaz provedeme indukcí.

    \begin{enumerate}
        \item Označíme si $A_h$ jako minimální počet vrcholů hloubkově vyváženého stromu hloubky h
        \item Ukážeme že počet vrcholů roste exponencionálně s hloubkou. Tedy hloubka roste logaritmicky díky inverzi těchto funkcí
        \item Pro h = 1 je strom s jedním vrcholem. Pro h = 2 je to strom s 2 vrcholy ...\\[-4\baselineskip]
        % pro tvar "isosceles triangle" a anchor .apex + výpočty souřadnic

\begin{figure}[H]
\centering
\begin{tikzpicture}[
    scale=1.35,
    every node/.style={
        circle,
        draw,
        fill=white,
        inner sep=0pt,
        minimum size=7pt
    },
    line width=0.7pt
]

% ============================================================
% GLOBÁLNÍ PARAMETRY
% ============================================================

% Vzdálenost sourozeneckých vrcholů
\def\SiblingDist{0.65}

% Vertikální vzdálenost mezi úrovněmi
\def\LevelDist{0.60}

\def\SubtreeLabelOffset{0.28} % svislá mezera mezi trojúhelníkem a popiskem


% ============================================================
% A_1
% ============================================================

\node (A1) at (0,0) {};

\node[draw=none, fill=none, circle] at (0,0.45)
    {$A_1=1$};


% ============================================================
% A_2
% ============================================================

\node (A2root) at (1.7,0) {};

\node (A2child) at
    ({1.7-\SiblingDist/2},-\LevelDist) {};

\draw (A2root) -- (A2child);

\node[draw=none, fill=none, circle] at (1.7,0.45)
    {$A_2=2$};


% ============================================================
% A_3
% ============================================================

\node (A3root) at (3.4,0) {};

\node (A3left) at
    ({3.4-\SiblingDist/2},-\LevelDist) {};

\node (A3right) at
    ({3.4+\SiblingDist/2},-\LevelDist) {};

\node (A3leftchild) at
    ({3.4-\SiblingDist},-{2*\LevelDist}) {};

\draw (A3root) -- (A3left);
\draw (A3root) -- (A3right);
\draw (A3left) -- (A3leftchild);

\node[draw=none, fill=none, circle] at (3.4,0.45)
    {$A_3=4$};


% ============================================================
% A_4
% ============================================================

\node (A4root) at (5.35,0) {};

% druhá úroveň
\node (A4left) at
    ({5.35-\SiblingDist/2},-\LevelDist) {};

\node (A4right) at
    ({5.35+\SiblingDist/2},-\LevelDist) {};

% třetí úroveň
\node (A4leftleft) at
    ({5.35-\SiblingDist},-{2*\LevelDist}) {};

\node (A4leftright) at
    (5.35,-{2*\LevelDist}) {};

\node (A4rightright) at
    ({5.35+\SiblingDist},-{2*\LevelDist}) {};

% čtvrtá úroveň
\node (A4bottom) at
    ({5.35-1.5*\SiblingDist},-{3*\LevelDist}) {};

% hrany
\draw (A4root) -- (A4left);
\draw (A4root) -- (A4right);

\draw (A4left) -- (A4leftleft);
\draw (A4left) -- (A4leftright);

\draw (A4right) -- (A4rightright);

\draw (A4leftleft) -- (A4bottom);

\node[draw=none, fill=none, circle] at (5.35,0.45)
    {$A_4=7$};


% ============================================================
% Rekurence A_h
% ============================================================

\node[
    draw=none,
    fill=none,
    circle,
    anchor=west
] at (6.55,0.45)
    {$A_h=A_{h-1}+A_{h-2}+1$};


% ============================================================
% kořen
% ============================================================

\node (Ahroot) at (7.9,0) {};


% ============================================================
% levý podstrom h-1
% ============================================================

\node[
    isosceles triangle,
    shape border rotate=90, % špička nahoru (místo doprava)
    anchor=apex,            % zarovnávej podle horní špičky
    draw,
    fill=white,
    minimum height=44pt,
    minimum width=26pt,
    inner sep=0pt
] (Ahleft) at
    ({7.9-\SiblingDist},-{1.25*\LevelDist})
    {};

% hrana do špičky trojúhelníku
\draw (Ahroot) -- (Ahleft.apex);


% ============================================================
% pravý podstrom h-2
% ============================================================

\node[
    isosceles triangle,
    shape border rotate=90, % špička nahoru (místo doprava)
    anchor=apex,            % zarovnávej podle horní špičky
    draw,
    fill=white,
    minimum height=26pt,
    minimum width=16pt,
    inner sep=0pt
] (Ahright) at
    ({7.9+\SiblingDist},-{1.25*\LevelDist})
    {};

% hrana do špičky trojúhelníku
\draw (Ahroot) -- (Ahright.apex);


% ============================================================
% popisky podstromů
% ============================================================

% umísti popisky vždy pod trojúhelník (nezávisle na jeho velikosti)
\node[draw=none, fill=none] at ($(Ahleft.south)+(0,-\SubtreeLabelOffset)$) {$h-1$};
\node[draw=none, fill=none] at ($(Ahright.south)+(0,-\SubtreeLabelOffset)$) {$h-2$};
    
\end{tikzpicture}
\caption{Rekurentní konstrukce minimálního AVL stromu $A_h = A_{h-1} + A_{h-2} + 1$.}
\end{figure}
        \item Obecně platí vztah $A_h = A_{h-1} + A_{h-2} + 1$
        \item Nyní dokážeme indukcí, že $A_{h+1} \geq 2^{h/2}$
        \item Základní krok $A_1 = 1 \geq 2^{0/2} = 1$ a $A_2 = 2 \geq 2^{1/2} = \sqrt{2}$
        \item Indukční krok:
        \begin{align*}
            A_{h+1} &= A_h + A_{h-1} + 1 \\
                    &> A_h + A_{h-1} \\
                    &= 2^{(h-1)/2} + 2^{(h-2)/2} \\
                    &= 2^{h/2}\cdot\bigl(2^{-1/2} + 2^{-1}\bigr) \\
                    &> 2^{h/2}\cdot\Bigl(\tfrac12 + \tfrac12\Bigr) \\
                    &= 2^{h/2}
        \end{align*}
        \item Hloubkově vyvážený strom hloubky h+1 má tedy nejméně ${2^{h/2}}$ vrcholů. Proto HV strom o $n$ vrcholech má hloubku nejvýše $\log_{\sqrt{2}}(n) + 1$

        \item Protože je binární, tak nemůže mít hloubu menší než $\floor{\log_{2}(n)}$
        \item Tedy hloubka HV stromu o n vrcholech je $\bigTheta{\log(n)}$
    \end{enumerate}
\end{proof}

\begin{remark}
    \noindent Detaily k důkazu:
    \begin{enumerate}
        \item V indukčním kroku využíváme $A_{h+1}$ místo $A_h$ je to kvůli případu kdy $h = 0$. V případě kdyby jsme použili $A_h$ museli by jsme explicitně vyřešit tento případ.
        \item V indukčním kroku využíváme exponencionálu $2^{h/2}$ ta se využívá jelikož to hezky vyjde. Ale je možné využít i jiný exponent. 
        \item Argument "protože je binární, tak nemůže mít hloubu menší než $\floor{\log_{2}(n)}$". Lze nahradit tím stejným důkazem. Pouze nyní s největším možným stromem, aby jsme ukázali, že platí $\bigO{\log_2{(n)}}$.
    \end{enumerate}
\end{remark} \bigskip

\subsection{Rotace AVL}

\begin{definition}
    AVL strom je hloubkově vyvážený BVS. Pro každý jeho vrchol $v$ definujeme \textbf{znaménko vrcholu} $\delta(v)$ jako rozdíl hloubek jeho pravého a levého podstromu:
    $$\delta(v) = h(R(v)) - h(L(v))$$
\end{definition} \medskip

\begin{definition}
V platném AVL stromu musí pro každý vrchol $v$ platit, že $\delta(v) \in \{-1, 0, +1\}$:
\begin{itemize}[itemsep=-0.2em]
    \item $\delta(v) = -1$: levý podstrom je o 1 hlubší než pravý,
    \item $\delta(v) = 0$: oba podstromy mají stejnou hloubku,
    \item $\delta(v) = +1$: pravý podstrom je o 1 hlubší než levý.
\end{itemize}    
\end{definition} \bigskip


\noindent V případě kdy $\delta(v) \notin \{-1, 0, 1\}$, pak musíme strom opravit pomocí jednoduchých či dvojitých rotací. Obecně rotace vypadá takto: \bigskip

% AVL rotace: pravá rotace (y,x,A,B,C) <-> levá rotace

\begin{figure}[H]
\centering
\begin{tikzpicture}[
    scale=1.15,
    every node/.style={
        circle,
        draw,
        fill=white,
        inner sep=0pt,
        minimum size=22pt,
        font=\Large\itshape
    },
    line width=0.9pt
]

% ============================================================
% GLOBÁLNÍ PARAMETRY
% ============================================================

\def\SiblingDist{1.15}  % vzdálenost sourozeneckých uzlů/trojúhelníků
\def\LevelDist{1.1}     % vertikální vzdálenost mezi úrovněmi
\def\SubtreeLabelOffset{0.30} % svislá mezera mezi trojúhelníkem a popiskem
\def\TriW{30pt}         % šířka trojúhelníku (podstromu)
\def\TriH{55pt}         % výška trojúhelníku (podstromu)

% ============================================================
% LEVÝ STROM  (kořen y, levý potomek x, pravý podstrom C)
% ============================================================

\node (Ly) at (0,0) {$y$};

\node (Lx) at (-\SiblingDist,-\LevelDist) {$x$};

\node[
    isosceles triangle,
    shape border rotate=90,
    anchor=apex,
    draw,
    fill=white,
    minimum height=\TriH,
    minimum width=\TriW,
    inner sep=0pt
] (LC) at ({0.9*\SiblingDist},-{1.15*\LevelDist}) {$C$};

\node[
    isosceles triangle,
    shape border rotate=90,
    anchor=apex,
    draw,
    fill=white,
    minimum height=\TriH,
    minimum width=\TriW,
    inner sep=0pt
] (LA) at ({-\SiblingDist-0.75*\SiblingDist},-{2.15*\LevelDist}) {$A$};

\node[
    isosceles triangle,
    shape border rotate=90,
    anchor=apex,
    draw,
    fill=white,
    minimum height=\TriH,
    minimum width=\TriW,
    inner sep=0pt
] (LB) at ({-\SiblingDist+0.75*\SiblingDist},-{2.15*\LevelDist}) {$B$};

\draw (Ly) -- (Lx);
\draw (Ly) -- (LC.apex);
\draw (Lx) -- (LA.apex);
\draw (Lx) -- (LB.apex);




% ============================================================
% PRAVÝ STROM  (kořen x, levý podstrom A, pravý potomek y)
% ============================================================

\def\XShift{6.5}

\node (Rx) at (\XShift,0) {$x$};

\node[
    isosceles triangle,
    shape border rotate=90,
    anchor=apex,
    draw,
    fill=white,
    minimum height=\TriH,
    minimum width=\TriW,
    inner sep=0pt
] (RA) at ({\XShift-0.9*\SiblingDist},-{1.15*\LevelDist}) {$A$};

\node (Ry) at ({\XShift+\SiblingDist},-\LevelDist) {$y$};

\node[
    isosceles triangle,
    shape border rotate=90,
    anchor=apex,
    draw,
    fill=white,
    minimum height=\TriH,
    minimum width=\TriW,
    inner sep=0pt
] (RB) at ({\XShift+\SiblingDist-0.75*\SiblingDist},-{2.15*\LevelDist}) {$B$};

\node[
    isosceles triangle,
    shape border rotate=90,
    anchor=apex,
    draw,
    fill=white,
    minimum height=\TriH,
    minimum width=\TriW,
    inner sep=0pt
] (RC) at ({\XShift+\SiblingDist+0.75*\SiblingDist},-{2.15*\LevelDist}) {$C$};

\draw (Rx) -- (RA.apex);
\draw (Rx) -- (Ry);
\draw (Ry) -- (RB.apex);
\draw (Ry) -- (RC.apex);




% ============================================================
% ŠIPKY ROTACÍ + POPISKY
% ============================================================

% velké šipky mezi kořeny (Ly = y vlevo) a (Rx = x vpravo)
% pomocné body pro stabilní centrování popisků
\path (Ly.east) ++(0,2mm) coordinate (RTopL);
\path (Rx.west) ++(0,2mm) coordinate (RTopR);
\path (Ly.east) ++(0,-2mm) coordinate (RBotL);
\path (Rx.west) ++(0,-2mm) coordinate (RBotR);

\draw[thick, -{Stealth[length=5mm,width=3.2mm]}]
    (RTopL) to[bend left=18] (RTopR);
\path (RTopL) -- (RTopR)
    node[midway, above=3mm,
         shape=rectangle, draw=none, fill=white,
         inner sep=1pt, font=\normalsize,
         text height=1.7ex, text depth=.3ex] {Pravá rotace};

\draw[thick, -{Stealth[length=5mm,width=3.2mm]}]
    (RBotR) to[bend left=18] (RBotL);
\path (RBotL) -- (RBotR)
    node[midway, below=3mm,
         shape=rectangle, draw=none, fill=white,
         inner sep=1pt, font=\normalsize,
         text height=1.7ex, text depth=.3ex] {Levá rotace};

\end{tikzpicture}
\caption{AVL rotace: pravá rotace (y,x,A,B,C) a levá rotace.}
\end{figure}
 \bigskip

\noindent Obrázek může prvně vypadat matoucí, ale je to kompletní návod jak je provést. Vrcholy $x$ a $y$ chceme pomocí rotace opravit. Trojuhelníky $A,B,C$ jsou jejich podstromy. \\

\noindent Rotací provádíme vždy když $|\delta(v)| \geq 2$. Kdy, ale máme zvolit levou a kdy pravou nebo dokonce dvojnou? Na tuto otázku si ihned odpovíme. \bigskip

%---------------------------------------------------------------
%-----------------------ROTACe----------------------------------
%---------------------------------------------------------------

\newcommand{\deltaVal}[1]{\delta(#1)}

\subsubsection{Kdy a jak se provádí rotace v AVL stromu}


Při provádění operací \texttt{AVLInsert} nebo \texttt{AVLDelete} se může hloubka některého podstromu změnit. Při návratu z rekurze směrem ke kořeni přepočítáme znaménka. Jakmile u nějakého vrcholu $x$ narazíme na hloubkovou nevyváženost, tj. $|\delta(x)| \ge 2$. Musíme tento vrchol a jeho podstrom opravit pomocí rotací. \medskip

\noindent Podle znaménka nevyváženého vrcholu $x$ a znaménka jeho příslušného syna $y$ rozlišujeme 4 typy rotací:

\begin{itemize}
    \item \textbf{Jednoduchá rotace doprava}: $\delta(x) = \mathbf{\textcolor{red}{-2}}$ a $\delta(y) = \mathbf{-1}$.
    \item \textbf{Jednoduchá rotace doleva}: $\delta(x) = \mathbf{\textcolor{red}{+2}}$ a $\delta(y) = \mathbf{+1}$.
    \item \textbf{Dvojitá rotace vlevo-vpravo}: $\delta(x) = \mathbf{\textcolor{red}{-2}}$ a $\delta(y) = \mathbf{+1}$.
    \item \textbf{Dvojitá rotace vpravo-vlevo}: $\delta(x) = \mathbf{\textcolor{red}{+2}}$ a $\delta(y) = \mathbf{-1}$.
\end{itemize}

\noindent Zjednodušeně řečeno pokud mají vrcholy $x,y$ stejné znaménko provádíme jednoduchou rotaci v opačném případě musíme provést dvojitou. \bigskip

% ------------------------------------------------------------ %
\subsubsection{Jednoduchá rotace doprava}
% ------------------------------------------------------------ %

\noindent U jednoduché rotace vpravo se připojí vrcholy $x$ jakožto pravý syn vrcholu $y$. Pouze strom B změní svou pozici na pravého syna vrcholu $x$.

\begin{figure}[H]
\centering
\begin{tikzpicture}[
  vnode/.style={circle, draw, minimum size=8mm, inner sep=0pt, font=\bfseries},
  stree_hplus/.style={isosceles triangle, draw, isosceles triangle apex angle=50, shape border rotate=90, minimum height=17mm, minimum width=13mm, inner sep=1pt, font=\small},
  stree_h/.style={isosceles triangle, draw, isosceles triangle apex angle=55, shape border rotate=90, minimum height=12mm, minimum width=10mm, inner sep=1pt, font=\small},
  stree_hminus/.style={isosceles triangle, draw, isosceles triangle apex angle=60, shape border rotate=90, minimum height=8mm, minimum width=8mm, inner sep=1pt, font=\scriptsize},
  thick, >=stealth,
  stree_h_custom/.style args={#1/#2/#3}{
    stree_h,
    isosceles triangle apex angle=#1,
    minimum height=#2,
    minimum width=#3
  },
]
  % STROM PŘED ROTACÍ
  \begin{scope}[shift={(-1.2,0)}]
    \node[vnode] (x) at (0, 1.8) {$x$};
    \node[red, right=1pt of x, font=\bfseries] {(\textcolor{red}{-2})};
    
    \node[vnode] (y) at (-1.5, 0.6) {$y$};
    \node[blue!80!black, left=1pt of y, font=\bfseries] {(-1)};
    
    % Podstrom A - h+1 (vysoký)
    \node[stree_h_custom=35/18mm/15mm, anchor=apex] (A) at (-2.3, -0.6) {$A$};
    \node[below=1pt of A, font=\scriptsize] {$h+1$};
    
    % Podstrom B - h (střední)
    \node[stree_h, anchor=apex] (B) at (-0.7, -0.6) {$B$};
    \node[below=1pt of B, font=\scriptsize] {$h$};
    
    \node[stree_h, anchor=apex] (C) at (1.2, 0.6) {$C$};
    \node[below=1pt of C, font=\scriptsize] {$h$};
    
    \draw (x) -- (y);
    \draw (y) -- (A.apex);
    \draw (y) -- (B.apex);
    \draw (x) -- (C.apex);
    
    \node at (0, -4) {\textbf{Strom před rotací}};
  \end{scope}

  % ŠIPKA ROTACE
  \draw[->, ultra thick, blue!70!black] (1.3, 0.6) -- node[above, font=\bfseries] {Rotace R} node[below, font=\scriptsize] {(doprava)} (4.2, 0.6);

  % STROM PO ROTACI
  \begin{scope}[shift={(7.5,0)}]
  
    \node[vnode] (y2) at (0, 1.8) {$y$};
    \node[green!50!black, right=1pt of y2, font=\bfseries] {(0)};
    
    % Podstrom A - h + 1 (vysoký)
    \node[stree_h_custom=35/18mm/15mm, anchor=apex] (A2) at (-1.5, 0.6) {$A$};
    \node[below=1pt of A2, font=\scriptsize] {$h + 1$};
    
    \node[vnode] (x2) at (1.3, 0.6) {$x$};
    \node[green!50!black, right=1pt of x2, font=\bfseries] {(0)};
    
    % Podstrom B - h
    \node[stree_h, anchor=apex] (B2) at (0.5, -0.6) {$B$};
    \node[below=1pt of B2, font=\scriptsize] {$h$};
    
    % Podstrom C - h
    \node[stree_h, anchor=apex] (C2) at (2.1, -0.6) {$C$};
    \node[below=1pt of C2, font=\scriptsize] {$h$};
    
    \draw (y2) -- (A2.apex);
    \draw (y2) -- (x2);
    \draw (x2) -- (B2.apex);
    \draw (x2) -- (C2.apex);
    
    \node at (0.4, -4) {\textbf{Strom po rotaci R}};
  \end{scope}
\end{tikzpicture}
\caption{Jednoduchá rotace doprava (R). Špatné znaménko $\mathbf{\textcolor{red}{-2}}$ u vrcholu $x$ je opraveno na $0$.}
\end{figure}

\bigskip

\subsubsection{Jednoduchá rotace doleva}
% ------------------------------------------------------------ %

\noindent U jednoduché rotace vlevo se připojí vrchol $x$ jako levý syn vrcholu $y$. Následně podstrom B se staně pravým synem vrcholu $x$.

\begin{figure}[H]
\centering
\begin{tikzpicture}[
  vnode/.style={circle, draw, minimum size=8mm, inner sep=0pt, font=\bfseries},
  stree_hplus/.style={isosceles triangle, draw, isosceles triangle apex angle=50, shape border rotate=90, minimum height=17mm, minimum width=13mm, inner sep=1pt, font=\small},
  stree_h/.style={isosceles triangle, draw, isosceles triangle apex angle=55, shape border rotate=90, minimum height=12mm, minimum width=10mm, inner sep=1pt, font=\small},
  stree_hminus/.style={isosceles triangle, draw, isosceles triangle apex angle=60, shape border rotate=90, minimum height=8mm, minimum width=8mm, inner sep=1pt, font=\scriptsize},
  thick, >=stealth,
  stree_h_custom/.style args={#1/#2/#3}{
    stree_h,
    isosceles triangle apex angle=#1,
    minimum height=#2,
    minimum width=#3
  },
]
  % STROM PŘED ROTACÍ
  \begin{scope}[shift={(-2,0)}]
    \node[vnode] (x) at (0, 1.8) {$x$};
    \node[red, left=1pt of x, font=\bfseries] {(\textcolor{red}{+2})};
    
    % Podstrom A - h (střední)
    \node[stree_h, anchor=apex] (A) at (-1.2, 0.6) {$A$};
    \node[below=1pt of A, font=\scriptsize] {$h$};
    
    \node[vnode] (y) at (1.5, 0.6) {$y$};
    \node[blue!80!black, left=1pt of y, font=\bfseries] {(+1)};
    
    % Podstrom B - h (střední)
    \node[stree_h, anchor=apex] (B) at (0.7, -0.6) {$B$};
    \node[below=1pt of B, font=\scriptsize] {$h$};
    
    % Podstrom C - h+1 (vysoký)
    \node[stree_h_custom=35/18mm/15mm, anchor=apex] (C) at (2.3, -0.6) {$C$};
    \node[below=1pt of C, font=\scriptsize] {$h+1$};
    
    \draw (x) -- (A.apex);
    \draw (x) -- (y);
    \draw (y) -- (B.apex);
    \draw (y) -- (C.apex);
    
    \node at (0.5, -4) {\textbf{Strom před rotací}};
  \end{scope}

  % ŠIPKA ROTACE
  \draw[->, ultra thick, blue!70!black] (1.8, 0.6) -- node[above, font=\bfseries] {Rotace L} node[below, font=\scriptsize] {(doleva)} (4.8, 0.6);

  % STROM PO ROTACI
  \begin{scope}[shift={(8.5,0)}]
    \node[vnode] (y2) at (0, 1.8) {$y$};
    \node[green!50!black, right=1pt of y2, font=\bfseries] {(0)};
    
    \node[vnode] (x2) at (-1.5, 0.6) {$x$};
    \node[green!50!black, right=1pt of x2, font=\bfseries] {(0)};
    
    % Podstrom A - h
    \node[stree_h, anchor=apex] (A2) at (-2.3, -0.6) {$A$};
    \node[below=1pt of A2, font=\scriptsize] {$h$};
    
    % Podstrom B - h
    \node[stree_h, anchor=apex] (B2) at (-0.7, -0.6) {$B$};
    \node[below=1pt of B2, font=\scriptsize] {$h$};
    
    % Podstrom C - h + 1 (vysoký)
    \node[stree_h_custom=35/18mm/15mm, anchor=apex] (C2) at (1.3, 0.6) {$C$};
    \node[below=1pt of C2, font=\scriptsize] {$h + 1$};
    
    \draw (y2) -- (x2);
    \draw (y2) -- (C2.apex);
    \draw (x2) -- (A2.apex);
    \draw (x2) -- (B2.apex);
    
    \node at (-0.4, -4) {\textbf{Strom po rotaci L}};
  \end{scope}
\end{tikzpicture}
\caption{Jednoduchá rotace doleva (L). Špatné znaménko $\mathbf{\textcolor{red}{+2}}$ u vrcholu $x$ je opraveno na $0$.}
\end{figure}

\newpage

\subsubsection{Dvojitá rotace vlevo-vpravo}
% ------------------------------------------------------------ %

Dvojitá rotace se provede jakožto dvě souvislé jednoduché rotace. První provedeme jednoduchou levou rotaci na vrchol $y$. Následně provedeme jednoduchou pravou rotaci na vrchol $x$.

\begin{figure}[H]
\centering
\begin{tikzpicture}[
  vnode/.style={circle, draw, minimum size=8mm, inner sep=0pt, font=\bfseries},
  stree_hplus/.style={isosceles triangle, draw, isosceles triangle apex angle=50, shape border rotate=90, minimum height=17mm, minimum width=13mm, inner sep=1pt, font=\small},
  stree_h/.style={isosceles triangle, draw, isosceles triangle apex angle=55, shape border rotate=90, minimum height=12mm, minimum width=10mm, inner sep=1pt, font=\small},
  stree_hminus/.style={isosceles triangle, draw, isosceles triangle apex angle=60, shape border rotate=90, minimum height=8mm, minimum width=8mm, inner sep=1pt, font=\scriptsize},
  thick, >=stealth,
  stree_h_custom/.style args={#1/#2/#3}{
    stree_h,
    isosceles triangle apex angle=#1,
    minimum height=#2,
    minimum width=#3
  },
]
% STROM PŘED ROTACÍ
\begin{scope}[shift={(-1,0)}]
    \node[vnode] (x) at (0, 2.4) {$x$};
    \node[red, right=1pt of x, font=\bfseries] {(\textcolor{red}{-2})};

    \node[vnode] (y) at (-1.9, 1.2) {$y$};
    \node[blue!80!black, left=1pt of y, font=\bfseries] {(+1)};

    % Podstrom A - h
    \node[stree_h, anchor=apex] (A) at (-3.0, 0.0) {$A$};
    \node[below=1pt of A, font=\scriptsize] {$h$};

    \node[vnode] (z) at (-0.8, 0.0) {$z$};
    \node[blue!80!black, right=1pt of z, font=\bfseries] {(-1)};

    % Podstrom B - h
    \node[stree_h, anchor=apex] (B) at (-1.5, -1.2) {$B$};
    \node[below=1pt of B, font=\scriptsize] {$h$};

    % Podstrom C - h-1
    \node[stree_hminus, anchor=apex] (C) at (0.1, -1.2) {$C$};
    \node[below=1pt of C, font=\scriptsize] {$h-1$};

    % Podstrom D - h
    \node[stree_h, anchor=apex] (D) at (1.5, 1.2) {$D$};
    \node[below=1pt of D, font=\scriptsize] {$h$};

    \draw (x) -- (y);
    \draw (x) -- (D.apex);
    \draw (y) -- (A.apex);
    \draw (y) -- (z);
    \draw (z) -- (B.apex);
    \draw (z) -- (C.apex);

    \node at (-0.5, -3.2) {\textbf{Strom před rotací LR}};
\end{scope}

  % ŠIPKA ROTACE
  \draw[->, ultra thick, blue!70!black] (1.9, 0.6) -- node[above, font=\bfseries] {Dvojitá LR} node[below, font=\scriptsize] {(L pak R)} (4.3, 0.6);

  % STROM PO ROTACI
\begin{scope}[shift={(8,-0.1)}]
    \node[vnode] (z2) at (0, 2.4) {$z$};
    \node[green!50!black, right=1pt of z2, font=\bfseries] {(0)};

    \node[vnode] (y2) at (-1.9, 1.2) {$y$};
    \node[green!50!black, right=1pt of y2, font=\bfseries] {(0)};

    % Podstrom A - h
    \node[stree_h, anchor=apex] (A2) at (-3.0, 0.0) {$A$};
    \node[below=1pt of A2, font=\scriptsize] {$h$};

    % Podstrom B - h
    \node[stree_h, anchor=apex] (B2) at (-1.4, 0.0) {$B$};
    \node[below=1pt of B2, font=\scriptsize] {$h$};

    \node[vnode] (x2) at (1.7, 1.2) {$x$};
    \node[blue!80!black, right=1pt of x2, font=\bfseries] {(+1)};

    % Podstrom C - h-1
    \node[stree_hminus, anchor=apex] (C2) at (0.6, 0.0) {$C$};
    \node[below=1pt of C2, font=\scriptsize] {$h-1$};

    % Podstrom D - h
    \node[stree_h, anchor=apex] (D2) at (2.8, 0.0) {$D$};
    \node[below=1pt of D2, font=\scriptsize] {$h$};

    \draw (z2) -- (y2);
    \draw (z2) -- (x2);
    \draw (y2) -- (A2.apex);
    \draw (y2) -- (B2.apex);
    \draw (x2) -- (C2.apex);
    \draw (x2) -- (D2.apex);

    \node at (0.0, -3.2) {\textbf{Strom po dvojité rotaci LR}};
\end{scope}
\end{tikzpicture}
\caption{Dvojitá rotace LR (pro případ $\delta(z) = -1$). Nevyváženost $\mathbf{\textcolor{red}{-2}}$ v kořeni $x$ je odstraněna překořeněním za vrchol $z$.}
\end{figure}

\bigskip

\subsubsection{Dvojitá rotace vpravo-vlevo}
% ------------------------------------------------------------ %

\noindent První provedeme jednoduchou levou pravou na vrchol $y$. Následně provedeme jednoduchou levou rotaci na vrchol $x$.

\begin{figure}[H]
\centering
\begin{tikzpicture}[
  vnode/.style={circle, draw, minimum size=8mm, inner sep=0pt, font=\bfseries},
  stree_hplus/.style={isosceles triangle, draw, isosceles triangle apex angle=50, shape border rotate=90, minimum height=17mm, minimum width=13mm, inner sep=1pt, font=\small},
  stree_h/.style={isosceles triangle, draw, isosceles triangle apex angle=55, shape border rotate=90, minimum height=12mm, minimum width=10mm, inner sep=1pt, font=\small},
  stree_hminus/.style={isosceles triangle, draw, isosceles triangle apex angle=60, shape border rotate=90, minimum height=8mm, minimum width=8mm, inner sep=1pt, font=\scriptsize},
  thick, >=stealth,
  stree_h_custom/.style args={#1/#2/#3}{
    stree_h,
    isosceles triangle apex angle=#1,
    minimum height=#2,
    minimum width=#3
  },
]
  % STROM PŘED ROTACÍ
  \begin{scope}[shift={(-1,0)}]
    \node[vnode] (x) at (0, 2.4) {$x$};
    \node[red, left=1pt of x, font=\bfseries] {(\textcolor{red}{+2})};

    % Podstrom A - h
    \node[stree_h, anchor=apex] (A) at (-1.4, 1.2) {$A$};
    \node[below=1pt of A, font=\scriptsize] {$h$};

    \node[vnode] (y) at (1.8, 1.2) {$y$};
    \node[blue!80!black, left=1pt of y, font=\bfseries] {(-1)};

    \node[vnode] (z) at (0.8, 0.0) {$z$};
    \node[blue!80!black, left=1pt of z, font=\bfseries] {(-1)};

    % Podstrom B - h
    \node[stree_h, anchor=apex] (B) at (0.1, -1.2) {$B$};
    \node[below=1pt of B, font=\scriptsize] {$h$};

    % Podstrom C - h-1
    \node[stree_hminus, anchor=apex] (C) at (1.4, -1.2) {$C$};
    \node[below=1pt of C, font=\scriptsize] {$h-1$};

    % Podstrom D - h
    \node[stree_h, anchor=apex] (D) at (2.9, 0.0) {$D$};
    \node[below=1pt of D, font=\scriptsize] {$h$};

    \draw (x) -- (A.apex);
    \draw (x) -- (y);
    \draw (y) -- (z);
    \draw (y) -- (D.apex);
    \draw (z) -- (B.apex);
    \draw (z) -- (C.apex);

    \node at (0.5, -3.2) {\textbf{Strom před rotací RL}};
  \end{scope}

  % ŠIPKA ROTACE
  \draw[->, ultra thick, blue!70!black] (3.1, 0.6) -- node[above, font=\bfseries] {Dvojitá RL} node[below, font=\scriptsize] {(R pak L)} (5.3, 0.6);

  % STROM PO ROTACI
  \begin{scope}[shift={(9,-0.1)}]
    \node[vnode] (z2) at (0, 2.4) {$z$};
    \node[green!50!black, right=1pt of z2, font=\bfseries] {(0)};

    \node[vnode] (x2) at (-1.8, 1.2) {$x$};
    \node[green!50!black, right=1pt of x2, font=\bfseries] {(0)};

    % Podstrom A - h
    \node[stree_h, anchor=apex] (A2) at (-2.6, 0.0) {$A$};
    \node[below=1pt of A2, font=\scriptsize] {$h$};

    % Podstrom B - h
    \node[stree_h, anchor=apex] (B2) at (-1.1, 0.0) {$B$};
    \node[below=1pt of B2, font=\scriptsize] {$h$};

    \node[vnode] (y2) at (1.6, 1.2) {$y$};
    \node[blue!80!black, right=1pt of y2, font=\bfseries] {(+1)};

    % Podstrom C - h-1
    \node[stree_hminus, anchor=apex] (C2) at (0.6, 0.0) {$C$};
    \node[below=1pt of C2, font=\scriptsize] {$h-1$};

    % Podstrom D - h
    \node[stree_h, anchor=apex] (D2) at (2.5, 0.0) {$D$};
    \node[below=1pt of D2, font=\scriptsize] {$h$};

    \draw (z2) -- (x2);
    \draw (z2) -- (y2);
    \draw (x2) -- (A2.apex);
    \draw (x2) -- (B2.apex);
    \draw (y2) -- (C2.apex);
    \draw (y2) -- (D2.apex);

    \node at (0.0, -3.2) {\textbf{Strom po dvojité rotaci RL}};
  \end{scope}
\end{tikzpicture}
\caption{Dvojitá rotace RL (pro případ $\delta(z) = -1$). Nevyváženost $\mathbf{\textcolor{red}{+2}}$ v kořeni $x$ je odstraněna překořeněním za vrchol $z$.}
\end{figure}


\newpage

% pseudokód
\subsubsection{Levá rotace pseudokód}

\noindent \textbf{Vstup}: Vrchol $x$ v BVS \newline
\funkce{RotateLeft}{vrchol $x$}
\begin{pseudocode}
    \item $y := r(x)$
    \item $r(x) := \ell(y)$
    \item $\ell(y) := x$
    \item $p(r(x)) := x;\ p(y) := p(x);\ p(x) := y;$
    \item $h(T(x)) := \max(h(L(x)),\, h(R(x))) + 1$
    \item $h(T(y)) := \max(h(R(y)),\, h(T(x))) + 1$
\end{pseudocode} \bigskip


% ============================================================ %
\subsection{Operace nad AVL Stromem}
% ============================================================ %

Jelikož je AVL strom je rozšířením BVS, všechny operace, které nemění strukturu ani obsah stromu, fungují v AVL stromech naprosto identicky jako v klasickém BVS. Operace, které měli časovou složitost, Jako BVSFind, BVSMin, BVSPred, atd. závislou na hloubce, nyní mají krásně logaritmickou složitost. \medskip

\noindent
Jak ale naložíme s operacemi které přidávají či mažou prvky? Ty mohou svým provedením porušit hloubkové vyvážení tedy $|\delta(v)| \geq 2$ pro nějaký vrchol $v$. Proto je musíme lehce modifikovat. První z nich si upravme přidávání.

% ------------------------------------------------------------ %
\subsubsection{Operace AVLInsert}
% ------------------------------------------------------------ %

\noindent Při vkládání budeme postupovat úplně stejně jako u BVSInsert. Při návratu z rekurze budeme muset přepočítat $\delta(v)$ pro všechny vrcholy $v$ od vloženého až do kořene. V případě pokud porušíme hloubkové vyvážení ($\delta(v) \geq 2$), poté musíme opravit strom pomocí rotací.

\bigskip

\funkce{AVLInsert}{ukazatel na vrchol $v$, klíč $x$}
\begin{pseudocode}
\item Pokud $v = \emptyset$:
\item \tab Vytvoř nový vrchol $v$ s klíčem $x$, $\delta(v) := 0$
\item \tab return $v$
\item Pokud $x = k(v)$: return $v$ \textcolor{gray}{// Klíč již ve stromu existuje}
\item Pokud $x < k(v)$:
\item \tab $\ell(v) := \text{\textcolor{blue}{AVLInsert}}(\ell(v), x)$
\item Pokud $x > k(v)$:
\item \tab $r(v) := \text{\textcolor{blue}{AVLInsert}}(r(v), x)$
\item $v := \text{\textcolor{blue}{VyvažAVL}}(v)$ \textcolor{gray}{// Přepočet znaménka $\delta(v)$ a případná rotace}
\item return $v$
\end{pseudocode}

\bigskip

\noindent Pokud porušíme hloubkovou vyváženost, tak po její opravě rotací se hloubka vrátí na původní hodnotu (jinak by stále pořušovala hloubkovou vyváženost). Tedy již nemusíme propagovat tuto změnu dále (z pohledu rodiče podstromu se nic nezměnilo).

% ------------------------------------------------------------ %
\subsubsection{Operace AVLDelete}
% ------------------------------------------------------------ %

Vrchol mažeme úplně stejně jako v BVS. Stejně tak při návratu z rekurze musíme přepočítat hodnotu $\delta(v)$, pro všechny $v$ na ceště od smazaného ke kořeni. A v případě $|\delta(v)| \geq 2$ opravit strom pomocí rotací. 

\bigskip
\funkce{AVLDelete}{ukazatel na vrchol $v$, klíč $x$}
\begin{pseudocode}
\item Pokud $v = \emptyset$: return $\emptyset$
\item Pokud $x < k(v)$: $\ell(v) := \text{\textcolor{blue}{AVLDelete}}(\ell(v), x)$
\item Pokud $x > k(v)$: $r(v) := \text{\textcolor{blue}{AVLDelete}}(r(v), x)$
\item Pokud $x = k(v)$:
\item \tab Pokud $\ell(v) = \emptyset$: return $r(v)$ \textcolor{gray}{// Nemá syna nebo má jen pravého}
\item \tab Pokud $r(v) = \emptyset$: return $\ell(v)$ \textcolor{gray}{// Má jen levého syna}
\item \tab $w := \text{\textcolor{blue}{BVSMin}}(r(v))$ \textcolor{gray}{// Má oba syny -> nahradíme následníkem}
\item \tab $k(v) := k(w)$
\item \tab $r(v) := \text{\textcolor{blue}{AVLDelete}}(r(v), k(w))$
\item $v := \text{\textcolor{blue}{VyvažAVL}}(v)$ \textcolor{gray}{// Přepočet znaménka $\delta(v)$ a případná rotace}
\item return $v$
\end{pseudocode}

\bigskip

\noindent Zde oproti AVLInsert nemůžeme propagaci zastavit po rotaci, jelikož se hloubka stromu mohla snížit o 1. Což dále může ovlivnit stupeň rodiče podstromu.

\subsection{Časová složitost operací nad AVL stromem}

\begin{theorem}
    AVLInsert a AVLDelete trvají $\bigO{log(n)}$.
\end{theorem}

\begin{proof}
    Důkaz provedeme přímo
    \begin{enumerate}
        \item Již víme že hloubka AVL stromu je vždy $\bigTheta{log(n)}$.
        \item Operace na BVS byli závislé na hloubce. Tedy určitě nebudou pracovat v lepším čase než logaritmickém tedy $\bigOmega{log(n)}$.
        \item Zároveň ale samotná rotace provede $\bigTheta{1}$ operací v každém vrcholu a to nejhůře na celé cestě do kořene. Což je celkem $\bigO{log(n)}$.
        \item Tedy celkem $\bigTheta{log(n)}$.
    \end{enumerate}
\end{proof}
